\section{State-and-Law Habitat}
\label{app:state-law-habitat}

This appendix records the compact convex state-and-law setting in which a
closure map of the observer layer would act.

The OPH inputs used here are:
\begin{enumerate}
\item the patch net \(P\mapsto A(P)\) of von Neumann algebras together with isotony and overlap
restriction maps;
\item a compatible local state family, plus the separately declared global-extension interface
used by this appendix to represent it on the inductive-limit algebra;
\item the finite premise-A3 information-projection branch and the separate optimizer-pushforward
interface, which together supply a common finite family of gauge-invariant local constraint
observables across cutoffs;
\item on the fixed-cutoff branch, the derived finite type-I presentations and compact boundary
gluing groups used elsewhere in the OPH package.
\end{enumerate}

The law slot is encoded by the finite expectation-value coordinates of the retained A1-generated
constraint family used by premise A3. That choice furnishes a compact convex state-and-law habitat
for the framework.

\paragraph{Theorem (Internal state-and-law habitat theorem).}
Fix finitely many screen patches \(P_1,\dots,P_N\) in the standard OPH setup. For each \(i\), let
\[
M_i:=A(P_i),
\qquad
M_{ij}:=A(P_i\cap P_j),
\]
where \(M_i\) is the patch von Neumann algebra and \(M_{ij}\subset M_i,M_j\) is the overlap
algebra. Let
\[
S_i:=S(M_i)\subset M_i^*
\]
be the full state space of \(M_i\), endowed with the weak* topology \(\sigma(M_i^*,M_i)\).

Choose a finite family of self-adjoint gauge-invariant local constraint observables
\[
\mathcal O=\{O_1,\dots,O_m\}
\]
from the A1-generated premise-A3 constraint grammar, with each \(O_a\) supported in one of the chosen patches; write
\(i(a)\) for an index such that \(O_a\in M_{i(a)}\).

Define the overlap-consistent state sector
\[
X_{\mathrm{ov}}
:=
\left\{
(\omega_1,\dots,\omega_N)\in \prod_{i=1}^N S_i:
\omega_i|_{M_{ij}}=\omega_j|_{M_{ij}}\ \text{for all }i,j
\right\}.
\]
Define the OPH law-coordinate map
\[
c:X_{\mathrm{ov}}\to \mathbb R^m,
\qquad
c(\omega_1,\dots,\omega_N)
:=
\bigl(\omega_{i(1)}(O_1),\dots,\omega_{i(m)}(O_m)\bigr).
\]
Let
\[
X_{\mathcal O}
:=
\left\{
(\omega_1,\dots,\omega_N,\ell)\in
\left(\prod_{i=1}^N S_i\right)\times \mathbb R^m:
(\omega_1,\dots,\omega_N)\in X_{\mathrm{ov}},
\ \ell=c(\omega_1,\dots,\omega_N)
\right\}.
\]
Then:
\begin{enumerate}
\item the ambient product
\[
E:=
\left(\prod_{i=1}^N M_i^*\right)\times \mathbb R^m
\]
is a Banach space for any product norm;
\item \(X_{\mathcal O}\) is nonempty and convex;
\item \(X_{\mathcal O}\) is norm-closed in \(E\);
\item \(X_{\mathcal O}\) is compact in the product topology
\[
\tau
:=
\left(\prod_{i=1}^N \sigma(M_i^*,M_i)\right)\times
\text{Euclidean topology on }\mathbb R^m;
\]
\end{enumerate}

\begin{proof}
Each \(M_i\) is a von Neumann algebra in the OPH patch net, hence a Banach space, so each dual
\(M_i^*\) is Banach. A finite product of Banach spaces is Banach, hence so is \(E\).

For each \(i\), the state space
\[
S_i=\{\omega\in M_i^*:\omega\ge 0,\ \omega(1)=1\}
\]
is convex and weak* compact by Banach--Alaoglu, because it is a weak* closed subset of the dual
unit ball. For every pair \((i,j)\), restriction along \(M_{ij}\subset M_i\) and
\(M_{ij}\subset M_j\) defines affine weak* continuous maps
\[
r_{ij}:S_i\to S(M_{ij}),
\qquad
r_{ji}:S_j\to S(M_{ij}).
\]
Hence \(X_{\mathrm{ov}}\) is an intersection of affine equalizer sets, so it is convex and
\(\tau\)-closed inside \(\prod_i S_i\). Since \(\prod_i S_i\) is compact and convex in the
product weak* topology, \(X_{\mathrm{ov}}\) is compact and convex as well.

Nonemptiness comes from the OPH global state on the inductive-limit algebra: its restrictions
\[
\bigl(\omega|_{M_1},\dots,\omega|_{M_N}\bigr)
\]
belong to \(X_{\mathrm{ov}}\) because the restrictions agree on every overlap by construction.

Each coordinate \(\omega_{i(a)}\mapsto \omega_{i(a)}(O_a)\) is affine and weak* continuous, since
evaluation at a fixed algebra element is weak* continuous. Thus \(c\) is affine and
\(\tau\)-continuous. The graph map
\[
G:X_{\mathrm{ov}}\to E,
\qquad
G(x):=(x,c(x)),
\]
is affine and continuous, so its image \(X_{\mathcal O}=G(X_{\mathrm{ov}})\) is convex and
\(\tau\)-compact.

To see norm-closedness in \(E\), let \((x_n,c(x_n))\to (x,\ell)\) in norm. Norm convergence
implies weak* convergence on each coordinate, so \(x\in X_{\mathrm{ov}}\) because
\(X_{\mathrm{ov}}\) is weak* closed. Since \(c\) is weak* continuous, \(c(x_n)\to c(x)\), while
norm convergence in \(\mathbb R^m\) also gives \(c(x_n)\to \ell\). Hence \(\ell=c(x)\), so
\((x,\ell)\in X_{\mathcal O}\).

Finally, \(X_{\mathcal O}\) is a compact convex subset of the locally convex topological vector
space \(E\) equipped with \(\tau\).
\end{proof}

The additional inputs needed for a strange-loop closure theorem are:
\begin{enumerate}
\item a closure map \(T\) built from internal operations such as overlap repair,
premise-A3 information reprojection, an explicit collar-recovery interface, and
record-sector coarse-graining;
\item a proof that the stronger ``observer-supporting'' or ``selected-world'' criteria carve out a
nonempty \(T\)-invariant observer-supporting subset of \(X_{\mathcal O}\);
\item any uniqueness, contraction, or Lyapunov-type stability estimate for that closure map.
\end{enumerate}

The habitat theorem supplies only the ambient Banach/compact-convex setting. It does not define a
canonical internal self-map, it does not identify any nonempty invariant observer-supporting subset,
and it does not prove uniqueness or stability. In particular, nonemptiness and compact-convexity
of the ambient habitat do not by themselves imply the existence of a nonempty invariant
observer-supporting sector. This does not erase the exact finite packet branch proved on the
consensus surface: there the quotient normal-form map pushes forward to a continuous
affine idempotent self-map of the finite packet simplex, with fixed points exactly the packets
supported on quotient normal forms. Any further strange-loop reading of the ambient habitat as a
closure story for existence is interpretive only and is not part of the recovered theorem package.

\section{Observer Registry and Clock-Naturality Certificate}
\label{app:observer-registry-clock-contract}

For finite exported observer systems, a theorem-faithful evidence bundle separates executor
metadata from observer semantics. It contains:
\begin{enumerate}[label=(\roman*),leftmargin=2em]
\item a semantic history DAG or labeled poset \(\mathsf H_O=(E_O,\preceq_O,\ell_O)\), with event
keys computed from semantic payload, observer token, visible footprint, and semantic parents;
\item a provenance envelope containing worker ID, repair iteration, retry count, queue position,
timestamp, and packet latency, none of which may enter the semantic event key unless declared as
physical input;
\item local observer registry groupoids with disjoint patch-observer, cap-observer, and
future-observer namespaces, explicit birth-event classes, continuation arrows, split/merge lineage
arrows, and overlap functors preserving those data;
\item a global registry descent certificate: cocycle agreement, trivial monodromy, no duplicate
global observer identifiers, no mixed patch/cap/future namespace collisions, and no reuse of local
anchor-patch identifiers as local observer indices;
\item observer-algebra extraction data
\[
\mathsf{ObsAlg}(\widehat q_{\mathrm{nf}})
=
(\mathcal A_O,\omega_O,\mathcal R_O,\ldots)
\]
whose outputs are state-preserving isomorphic for history-equivalent or implementation-equivalent
terminal augmented normal forms;
\item a support-cap chart map \(J_O\) natural with those observer-algebra isomorphisms; and
\item a clock instrument \(\Theta_O\) with affine reparameterization and residual bound for state,
calibration, and record-readout errors.
\end{enumerate}

\begin{theorem}[Finite observer-clock certificate soundness]\label{thm:finite-observer-clock-certificate}
On a finite exported model, suppose the augmented repair relation on
\(\widehat Q=\{(q,[\mathsf H])\}\) terminates, terminal states are exactly the consistent physical
states with completed semantic records, every augmented critical pair has a history-coherent join,
implementation steps are linearizable to semantic commits or stutters, registry descent passes the
groupoid cocycle and monodromy checks, and the observer-algebra/chart/clock maps satisfy the
state-preserving and affine-residual certificates above. Then worker count, queue order, retry
delivery, scheduler counters, and packet latency that is not declared physical input do not change
the observer-readable history class, registry class, or clock law beyond the declared affine clock
residual.
\end{theorem}

\begin{proof}
Termination plus the augmented local diamond gives confluence by Newman's lemma for terminal
history classes rather than only physical quotient states. Linearizability says every
implementation execution projects to the same semantic relation up to stutters, so executor
metadata cannot create new semantic events. Registry descent identifies the same global observer
groupoid across local presentations. State-preserving observer-algebra extraction and natural
support-cap charts transport the modular flow along the same observer identity, while the clock
instrument certificate bounds the operational readout after the declared affine reparameterization.
\end{proof}

\begin{remark}[Falsifiers]\label{rem:observer-clock-falsifiers}
The certificate fails if semantic event identifiers depend on worker IDs, timestamps, queue
positions, repair iteration counters, or retries; if duplicate retry delivery creates two visible
record commits; if two namespaces collide in the global registry; if registry cocycles or lineage
arrows fail on overlaps; if histories from two executors are not labeled-poset isomorphic; or if
clock residuals exceed the declared bound. These are public evidence-bundle failures that exceed
field-name mismatches.
\end{remark}

